\usepackage[utf8]{inputenc}
\usepackage[T1]{fontenc}
\usepackage[ngerman]{babel}
\usepackage{array}
\usepackage{bookmark}
\usepackage{textpos}
\usepackage{wasysym}
\usepackage{listliketab}
\usepackage[NewCommands,NewParameters]{ragged2e}
\usepackage{tikz}
\usepackage{eurosym}
\usepackage{textpos}
\usepackage{comment}

\usetheme{junghans}

\def\insertauthorindicator{Wer?}% Default is "Who?"
\def\insertinstituteindicator{Von?}% Default is "From?"
\def\insertdateindicator{Wann?}% Default is "When?"

%\setbeamertemplate{caption}[numbered]

%%neuer Spaltentyp: linksb�ndig + Silbentrennung
\newcolumntype{M}[1]{>{\RaggedRight\hspace{0pt}\arraybackslash}p{#1}}

\author{Marius Spix}

\institute[\raisebox{-1ex}{\includegraphics[height=0.24cm]{images/JW_Kontur}}]{\includegraphics[scale=0.15]{images/J_60_Schwarz_Reflex_Blue}}
\title{QM-Cockpit}

\addtobeamertemplate{title page}{\vspace{1\baselineskip}\hspace{-9ex}\includegraphics[scale=.5]{images/Flaticon_SVGFiles/platform11}\vspace{-2\baselineskip}}

%\addtobeamertemplate{frametitle}{}{
%\begin{textblock*}{100mm}(-4.15cm,-.55cm)
%\includegraphics[scale=.33]{images/Fahne2}\end{textblock*}
%}

\AtBeginSection{\frame{\sectionpage}}
\AtBeginSubsection{\frame{\subsectionpage}}
\newtranslation[to=ngerman]{Section}{Abschnitt}
\newtranslation[to=ngerman]{Subsection}{Unterabschnitt}

\keywords{Abschlusspr�fung, Fachinformatiker, SAP, Quali"atsmanagement,
Lager, Warenwirtschaft}
\date{\today}

\frenchspacing

\begin{document}

\setcounter{framenumber}{0} 
\begin{frame}
  \maketitle
\end{frame}

\setcounter{figure}{0}

\begin{frame}{"Ubersicht}
  \tableofcontents
\end{frame}

\newcommand{\ipsum}[1]{
\subsection{{#1}}
\begin{frame}{Lorem Ipsum}
\onslide<2->{
\begin{block}{Lorem}
\begin{itemize}
  \item<3-> Lorem ipsum
  \item<4-> Dolor est
  \item<5-> Sit amit
  \item<6-> Bla
\end{itemize}
\end{block}
}
\onslide<8->{
\begin{block}{Ipsum}
\begin{enumerate}
  \item<9-> eins
  \item<10-> zwei
  \item<11-> drei
\end{enumerate}
\end{block}
}
\end{frame}
}


\section{Projektumfeld}
%Folie Junghans
\begin{frame}{Das Unternehmen}
	\begin{block}{foo}
	\end{block}

	\begin{block}{Plattform}
		\begin{itemize}
			\item{seit Sommer 2011}
			\item{SAP ECC 6.1}
		\end{itemize}
	\end{block}
\end{frame}

%Folie Prozesse im Wareneingang


\begin{comment}
\subsection{Das Unternehmen}
\begin{frame}{Junghans Unternehmensgruppe}
			\begin{block}{Fakten}
				\begin{itemize}
						\item Versandhandelsgruppe
						\item Hauptsitz in Aachen
						\item ca. 540 Mitarbeiter
						\item vertreten in 7 L"andern
				\end{itemize}
			\end{block}
			\begin{block}{Fakten}
				\begin{itemize}
						\item yyyy aktive Kunden
						\item \EUR{137\,000\,000} Umsatz %Gesch�ftsjahr 2012/13
						\item 1\,000\,000 Auftr�ge
				\end{itemize}
			\end{block}
\end{frame}

\begin{frame}{Produkte}
	\begin{block}{Junghans Wolle}
		\begin{itemize}
			\item Wolle und Garne
			\item Handarbeitszubeh"or
			\item K"unstlerbedarf
		\end{itemize}
	\end{block}
	\begin{block}{Pro-Idee}
		\begin{itemize}
			\item innovative Produkte aus aller Welt
			\item Bekleidung
			\item internationale Weine
			\item Kunstwerke
			\item M"obel, K"uchenutensilien
		\end{itemize}
	\end{block}
\end{frame}


\subsection{Die DV}
\begin{frame}{Organisation \& Datenverarbeitung}
baz

%Mitarbeiter
\end{frame}

\begin{frame}{SAP-System}
	\begin{block}{Allgemeines}
		\begin{itemize}
			\item seit Sommer 2011
			\item E
			\item
		\end{itemize}
	\end{block}
	\begin{block}{Drei Systeme}
		\begin{itemize}
			\item seit Sommer 2011
			\item E
			\item
		\end{itemize}
	\end{block}
\end{frame}
\end{comment}


\section{Ist-Analyse}
\begin{frame}{Analyse}
	\begin{block}{Methoden}
		\begin{itemize}
			\item Arbeitsbeobachtung
			\item Interview
		\end{itemize}
	\end{block}
	\begin{block}{Erkenntnisse}
		\begin{itemize}
			\item saisonale Schwankungen
			\item Personalplanung
			\item keine Standardfunktionalit"at
		\end{itemize}
	\end{block}
\end{frame}

\begin{frame}{Die Folge}

\begin{columns}
       \transboxout<2->
\begin{column}{.45\textwidth}
	\onslide<2->{\includegraphics[height=.42\textheight]{images/metro_ui_icon_set/Folders_OS/Documents.png}}\\
	\onslide<4->{\includegraphics[height=.42\textheight]{images/metro_ui_icon_set/Office_Apps/Word_2013.png}}
\end{column}
\begin{column}{.45\textwidth}
	\onslide<3->{\includegraphics[height=.42\textheight]{images/metro_ui_icon_set/Office_Apps/Excel_2013.png}}\\
	\onslide<5->{\includegraphics[height=.42\textheight]{images/metro_ui_icon_set/Other/Mail.png}}
\end{column}
\end{columns}

%Hier Grafik mit Unordnung
\end{frame}




\section{Soll-Konzept}


\section{Durchf"uhrung}
\begin{frame}{Entwicklungsprozess}
\end{frame}

\section{Ergebnis}
\ipsum{Bla}

{ %schwarze Folie
\setbeamercolor{background canvas}{bg=black}
\frame[plain]{}
} 

\end{document}

